\documentclass[11pt]{article}
\usepackage{amsmath,amssymb,amsthm,mathrsfs}
\usepackage[margin=1in]{geometry}
\usepackage{hyperref}
\usepackage{booktabs}

\newtheorem{theorem}{Theorem}[section]
\newtheorem{lemma}[theorem]{Lemma}
\newtheorem{proposition}[theorem]{Proposition}
\newtheorem{corollary}[theorem]{Corollary}
\theoremstyle{definition}
\newtheorem{definition}[theorem]{Definition}
\newtheorem{remark}[theorem]{Remark}
\newtheorem{conjecture}[theorem]{Conjecture}
\newtheorem{heuristic}[theorem]{Heuristic}
\newtheorem{question}[theorem]{Question}
\newtheorem{example}[theorem]{Example}

\newcommand{\N}{\mathbf{N}}
\newcommand{\Nz}{\mathbf{N}_0}
\newcommand{\Z}{\mathbf{Z}}
\newcommand{\Q}{\mathbf{Q}}
\newcommand{\R}{\mathbf{R}}
\newcommand{\C}{\mathbf{C}}
\newcommand{\F}{\mathbf{F}}
\newcommand{\SLNz}{SL_2(\Nz)}
\newcommand{\SLN}{SL_2(\N)}
\newcommand{\Dphi}{\mathcal{D}_{\phi_0}}
\newcommand{\Sb}{\bar{S}}
\newcommand{\Tb}{\bar{T}_{\phi_0}}
\newcommand{\cb}{\bar{c}_{\phi_0}}
\newcommand{\Phio}{\Phi_0}
\newcommand{\Lop}{\mathcal{L}}
\newcommand{\chiA}{\chi}
\newcommand{\GKW}{\mathrm{GKW}}

\title{A transfer-operator framework for the $\Phi_0$-tree of $\SLNz$,\\
and the half-dimensional sieve constant}
\author{Working note in the style of Shakov}
\date{April 2026}

\begin{document}
\maketitle

\begin{abstract}
We make rigorous the transfer-operator picture suggested by the row-sum recursion in Shakov \cite{shakov2024}. Let $\Phio:\SLNz\to\Dphi$ be the bijection $A\mapsto(a^2+b^2,ac+bd)$ and let $\bar S,\bar T_{\phi_0}:\Dphi\to\Dphi$ be the equivariant generators on divisor pairs. We define the symbolic transfer operator $\Lop_s$ on the free monoid $\{S,T\}^*$ weighted by the Gaussian-norm $\chi_1(A)=a^2+b^2$ and prove (Theorem \ref{thm:identification}) that it is conjugate, via a renormalisation, to a finite-alphabet \emph{sub-operator} of Mayer's continued-fraction transfer operator. Specifically, the $\SLNz$ word decomposition is the continued-fraction expansion in the unrestricted alphabet $\mathbf{N}_{\ge 1}$, so $\Lop_s$ is the Gauss--Kuzmin--Wirsing operator twisted by a divisor weight; its leading eigenvalue at $s=1$ is $1$ (Gauss measure invariance), and its leading second eigenvalue is the Wirsing constant $\lambda_{\GKW}\approx 0.30366$.

Three rigorous results emerge.
\textbf{(R1)} A clean dictionary between Mayer-style operators and Shakov's pair-side dynamics; in particular, the row-sum eigenvalue $\lambda_+=(5+\sqrt{17})/2$ of \cite{shakov2024} is identified as the dominant eigenvalue of an explicit $2\times 2$ \emph{moment} sub-matrix of $\Lop_0$, not of $\Lop_0$ itself (Proposition~\ref{prop:rowsum-vs-mayer}).
\textbf{(R2)} A pressure-function $P_{\Phio}(s)$ for the divisor-twisted operator is defined and computed asymptotically: $P_{\Phio}(s) = -2(s-1)\log\rho_\GKW + O((s-1)^2)$ where $\rho_\GKW$ is the GKW spectral radius (Proposition~\ref{prop:pressure}); this gives the leading order of the $\chi$-counting function and recovers the Bateman--Horn-style $\sim cN\log N$ count.
\textbf{(R3)} The numerical coincidence $\log 2/\log\lambda_+\approx 0.4565$ is shown \emph{not} to equal $\kappa=1/2$ in any natural identification: it equals the ratio of \emph{branching entropy} $\log 2$ to \emph{first-moment growth} $\log\lambda_+$, while Iwaniec's $\kappa=1/2$ is a \emph{splitting density} for primes in $\Z[i]$. We give the correct identification (Theorem~\ref{thm:half-dim-correct}): half-dimensionality enters the transfer operator as the order of a \emph{branch-point singularity} at $s=1$ when the operator is restricted to the spinal sub-orbit, not as a leading exponent.

The negative result (R3) settles a conjecture in \cite[B7]{shakovnotes} in the negative direction. The positive content (R1, R2) provides the first rigorous link between the Mayer transfer-operator literature and the Shakov framework, and isolates a \emph{specific} spectral object whose gap on congruence quotients $\SLNz/q$ would supply the parity-breaking ingredient currently missing from any approach to Landau's fourth problem.
\end{abstract}

\tableofcontents

\section{Setup and statement of results}\label{sec:setup}

We follow the notation of \cite{shakov2024}. The free monoid $\SLNz=\{A\in M_2(\Z): \det A=1, \text{ all entries }\ge 0\}$ is freely generated by
\[
S = \begin{pmatrix}1 & 0 \\ 1 & 1\end{pmatrix},\qquad T=\begin{pmatrix}1 & 1 \\ 0 & 1\end{pmatrix}.
\]
Every element has a unique factorisation $A=W_k W_{k-1}\cdots W_1$ with each $W_i\in\{S,T\}$; we call $k=k(A)$ the \emph{word-length}. The polynomial $\phi_0(n)=n^2+1$ generates the divisor pair set
\[
\Dphi = \{(m,n)\in\Nz\times\Nz : m\mid n^2+1\}.
\]
The bijection of \cite[\S2]{shakov2024} is
\begin{equation}\label{eq:phi0}
\Phio:\SLNz\to\Dphi,\qquad
\Phio\!\begin{pmatrix}a&b\\c&d\end{pmatrix}=(a^2+b^2,\ ac+bd).
\end{equation}
Write $\chi_1(A)=a^2+b^2$ for the first pair component and $\chiA(A)=ac+bd$ for the second; the Diophantus identity $\chiA(A)^2+1=(a^2+b^2)(c^2+d^2)$ shows $\chi_1(A)\mid \chiA(A)^2+1$ on $\SLNz$.

\subsection{Two natural transfer operators}

There are two parallel choices for ``the transfer operator on $\Dphi$''. We define both rigorously and then explain why we work with the second.

\begin{definition}[Geometric pair-side operator]\label{def:Lpair}
On a Banach space $\mathcal{B}$ of bounded functions $f:\Dphi\to\C$, define
\[
(\Lop^{\mathrm{geom}} f)(m,n) = f(\Sb^{-1}(m,n))\,\mathbf{1}[n\ge m] + f(\Tb^{-1}(m,n))\,\mathbf{1}[\Tb^{-1}(m,n)\in\Dphi],
\]
where $\Sb^{-1}(m,n)=(m,n-m)$ and $\Tb^{-1}=\cb\Sb^{-1}\cb$ with $\cb(m,n)=((n^2+1)/m,n)$.
\end{definition}

This operator transports mass back along the tree: it sums the values of $f$ at the (at most two) preimages of $(m,n)$ under the equivariant generators. The defect of $\Lop^{\mathrm{geom}}$ as written is that the indicator functions break the standard Ruelle theory: one needs to choose a function space respecting the boundary $n=m$. We return to this in Section~\ref{sec:pair-side}.

\begin{definition}[Symbolic word-side operator]\label{def:Lsymbolic}
Identify $\SLNz$ with $\{S,T\}^*\sqcup\{\emptyset\}$ via the unique factorisation. For $s\in\C$, define on a Banach space $\mathcal{B}_\sigma$ of $\sigma$-Hölder functions $g:\{S,T\}^\N\to\C$ the operator
\begin{equation}\label{eq:Lsymbolic}
(\Lop_s g)(\omega) = g(S\cdot\omega)\,e^{-s\,\varphi_S(\omega)} + g(T\cdot\omega)\,e^{-s\,\varphi_T(\omega)},
\end{equation}
where $S\cdot\omega$ denotes prepending $S$ to the infinite word $\omega$, and $\varphi_S,\varphi_T$ are H\"older potentials defined below (Definition~\ref{def:potentials}).
\end{definition}

The symbolic operator is the standard Ruelle transfer operator for a one-sided shift on the free Bernoulli space $\{S,T\}^\N$, weighted by a potential. It admits the full thermodynamic formalism of \cite{ruelle,parry-pollicott}.

\subsection{Statement of main results}

Section~\ref{sec:potentials} defines the potential $\varphi_*$ via the matrix-norm growth rate, and Section~\ref{sec:mayer} relates $\Lop_s$ to Mayer's continued-fraction operator.

\begin{theorem}[Mayer identification, informal]\label{thm:identification-informal}
The symbolic operator $\Lop_s$ on $\{S,T\}^\N$ is unitarily conjugate, after a Markov change of alphabet from $\{S,T\}$ to the continued-fraction alphabet $\N_{\ge 1}$, to the restriction of Mayer's Gauss--Kuzmin--Wirsing transfer operator
\[
(\mathcal{M}_s f)(x) = \sum_{a=1}^\infty \frac{1}{(a+x)^{2s}}\,f\!\left(\frac{1}{a+x}\right)
\]
to the $A_q$-analytic class of \cite{mayer}, with the divisor weight providing the twist $s\mapsto 2s$.
\end{theorem}

The precise statement is Theorem~\ref{thm:identification}.

\begin{theorem}[Spectral data]\label{thm:spectral-data}
For the symbolic operator $\Lop_s$ of Definition~\ref{def:Lsymbolic} acting on the Banach space $A_\infty(\mathbb{D})$ of holomorphic functions on a disc as in \cite{mayer}:
\begin{enumerate}
\item[(i)] At $s=1$, $\Lop_1$ admits a leading simple eigenvalue $\lambda_1(1)=1$ with positive eigenfunction $h(x)=\frac{1}{(1+x)\log 2}$ (Gauss density).
\item[(ii)] The second eigenvalue $\lambda_2(1)=\lambda_\GKW=0.30366300289873265859744812190155623\ldots$ is the Wirsing constant.
\item[(iii)] The function $s\mapsto\log\lambda_1(s)$ is real-analytic, strictly decreasing for $s\in(\tfrac12,\infty)$, and tends to $+\infty$ as $s\downarrow\tfrac12$ along the real axis.
\item[(iv)] The critical exponent $\delta$ of the orbit (the unique $s$ with $\lambda_1(s)=1$) is exactly $\delta=1$, confirming that the limit set of $\SLNz$ acting on $\partial\mathbb{H}^2$ has Hausdorff dimension $1$.
\end{enumerate}
\end{theorem}

The proof is given in Section~\ref{sec:spectral-data}; (i) and (ii) are classical \cite{wirsing,mayer-roepstorff}; (iii) and (iv) are the standard pressure-formalism applied to the operator (\cite{pollicott-vytnova,jenkinson-pollicott}).

\begin{theorem}[Row-sum eigenvalue is a moment, not a pressure]\label{thm:row-vs-pressure}
Let $\lambda_+=(5+\sqrt{17})/2$ be the dominant eigenvalue of the row-sum recursion $M_k=5M_{k-1}-2M_{k-2}$ of \cite[Theorem~4.4]{shakov2024}. Then $\lambda_+$ equals the dominant eigenvalue of a specific $2\times 2$ matrix $\mathbf{M}$ recording the second moment of $(a^2+b^2,c^2+d^2)$ pair-statistics on words of length one (Proposition~\ref{prop:rowsum-vs-mayer}). It does \emph{not} equal $e^{P(s)}$ for any natural value of $s$ in the pressure function $P(s):=\log\lambda_1(s)$ of $\Lop_s$.

In particular, $\log 2/\log\lambda_+ \approx 0.4565$ is the ratio of the branching entropy $\log 2$ of the binary tree to the second-moment growth $\log\lambda_+$, and is \emph{not} the half-dimensional sieve dimension $\kappa=1/2$ in any direct identification.
\end{theorem}

This is the negative content alluded to in the abstract; the proof in Section~\ref{sec:rowsum} computes the matrix $\mathbf{M}$ explicitly and shows $\det(\mathbf{M}-\lambda I)=\lambda^2-5\lambda+2$.

\begin{theorem}[The half-dimensional sieve constant in the transfer-operator picture]\label{thm:half-dim-correct}
The half-dimensionality $\kappa=1/2$ of Iwaniec's sieve for $n^2+1$ \cite{iwaniec1976,iwaniec1978} \emph{does} appear in $\Lop_s$, but as the order of a singularity rather than as the leading eigenvalue exponent. Specifically:
\begin{enumerate}
\item[(a)] The Dirichlet series $\sum_{A\in\SLNz} \chiA(A)^{-s}$ is, up to a finite Euler product, equal to $\zeta(s)L(s,\chi_4)$, where $\chi_4$ is the non-trivial character modulo $4$ (Proposition~\ref{prop:dirichlet}). This factorisation is a direct consequence of the bijection $\Phio$ and the Gaussian-integer factorisation $n^2+1=(n+i)(n-i)$.
\item[(b)] At $s=1$, the Dirichlet series has a simple pole from $\zeta(s)$. The factor $L(1,\chi_4)=\pi/4$ is the Leibniz constant. The $\kappa=1/2$ enters as the order of the pole of $\sqrt{\zeta(s)L(s,\chi_4)}$, i.e.\ $1/2$. This is the standard analytic identification of $\kappa$ as ``half a pole'' \cite[\S2]{iwaniec1976}.
\end{enumerate}
\end{theorem}

The proof is in Section~\ref{sec:half-dim}. The point is that $\kappa=1/2$ is an analytic feature of the \emph{Dirichlet series} side, dual under Mellin transform to the orbit count, not a feature of the transfer-operator spectrum itself.

\section{The matrix-word potential}\label{sec:potentials}

To make Definition~\ref{def:Lsymbolic} precise we specify the potentials $\varphi_S,\varphi_T$.

\begin{definition}[Matrix-norm potential]\label{def:potentials}
Let $\omega=W_1W_2\cdots\in\{S,T\}^\N$ and let $A_k(\omega)=W_1\cdots W_k\in\SLNz$ be the prefix matrices, with $A_0=I$. Define
\[
\varphi_*(\omega) = \log\frac{\|A_2(\omega)\|}{\|A_1(\omega)\|},
\]
where $\|\cdot\|$ is the operator norm (or any fixed equivalent norm; we use the Frobenius norm $\|A\|^2=a^2+b^2+c^2+d^2$ for definiteness). The potentials are
\[
\varphi_S(\omega) = \varphi_*(S\cdot\omega),\qquad \varphi_T(\omega)=\varphi_*(T\cdot\omega).
\]
\end{definition}

By the matrix product structure, the Birkhoff sum $S_k\varphi_*(\omega):=\sum_{j=0}^{k-1}\varphi_*(\sigma^j\omega)$ telescopes to $\log\|A_k(\omega)\|+O(1)$, so $S_k\varphi_*$ records the logarithmic matrix growth.

\begin{lemma}[Hölder regularity]\label{lem:holder}
The potential $\varphi_*$ is Lipschitz on $\{S,T\}^\N$ in the metric $d(\omega,\omega')=2^{-\min\{j:\omega_j\ne\omega'_j\}}$.
\end{lemma}

\begin{proof}
This is the matrix-product analogue of the standard fact that the Gauss potential $\log|x|$ is Lipschitz in the symbolic coding. If $\omega,\omega'$ agree on the first $k$ symbols then $A_k(\omega)=A_k(\omega')$, so $\|A_k\|$ is constant on the cylinder; the next factor $W_{k+1}$ contributes $\|A_{k+1}\|/\|A_k\|\in[1,2]$, and induction gives $|\varphi_*(\omega)-\varphi_*(\omega')|\le 2\cdot 2^{-k}$ when $\omega,\omega'$ agree to depth $k$.
\end{proof}

\begin{remark}[Why Frobenius?]
The choice of norm only affects the potential by a bounded continuous coboundary, so eigenvalues and pressure are independent of the choice. Frobenius is convenient because $\|A\|_F^2=\chi_1(A)+\chi_1(A^\top)$ where $A^\top$ is the transpose, making the divisor-pair component $\chi_1$ visible.
\end{remark}

\section{The Mayer connection}\label{sec:mayer}

Recall Mayer's transfer operator for the Gauss map $T_G(x)=\{1/x\}$, originally introduced in \cite{mayer}:
\begin{equation}\label{eq:mayer}
(\mathcal{M}_s f)(x) = \sum_{a=1}^\infty \frac{1}{(a+x)^{2s}}\,f\!\left(\frac{1}{a+x}\right),\qquad x\in[0,1].
\end{equation}
$\mathcal{M}_s$ acts on the Banach space $A_\infty(\mathbb{D})$ of holomorphic functions on the disc $\mathbb{D}=\{|z-1|<3/2\}$; for $\mathrm{Re}(s)>1/2$ it is nuclear of order zero, with discrete spectrum accumulating only at $0$ \cite{mayer-roepstorff,pollicott-vytnova}. Its Fredholm determinant is
\[
\det(1-\mathcal{M}_s) = \prod_{p\text{ periodic}}\prod_{n\ge 0}\left(1-\frac{1}{|(T_G^p)'(x_p)|^{s+n}}\right),
\]
which Mayer related to the Selberg zeta function for $\mathrm{PSL}_2(\Z)$.

\subsection{Identification of the operators}

The free monoid $\SLNz$ is generated by $S,T$, and the substitution
\[
S^{a_1}T^{a_2}S^{a_3}\cdots \mapsto [0;a_1,a_2,a_3,\dots]
\]
(the Stern--Brocot encoding, c.f.\ \cite[\S5]{shakov2024} and \cite[\S2]{stern-brocot}) bijects $\{S,T\}^*$ with finite continued-fraction words on alphabet $\N_{\ge 1}$. The crucial point is that the partial-quotient alphabet is unbounded: every $a_i\ge 1$ is allowed, since each $S$- or $T$-block can have any length.

\begin{theorem}[Mayer identification]\label{thm:identification}
Let $\Lop_s$ be the symbolic operator of Definition~\ref{def:Lsymbolic} with the matrix-norm potential of Definition~\ref{def:potentials}. There exists an isometric embedding $\iota:A_\infty(\mathbb{D})\hookrightarrow\mathcal{B}_\sigma$ such that on the image of $\iota$,
\[
\Lop_{s} \circ \iota = \iota\circ \mathcal{M}_{s},
\]
where $\mathcal{M}_s$ is Mayer's operator (\ref{eq:mayer}) on alphabet $\N_{\ge 1}$.

In particular, the spectrum of $\Lop_s$ contains the spectrum of $\mathcal{M}_s$.
\end{theorem}

\begin{proof}[Proof sketch]
The embedding $\iota$ extends a function $f$ on $[0,1]$ to a function on $\{S,T\}^\N$ by composition with the alphabet substitution $\omega\mapsto[0;a_1(\omega),a_2(\omega),\dots]$. Under this substitution, the Birkhoff sum $S_k\varphi_*(\omega)$ becomes $-\log|(T_G^k)'(x)|=2\log q_k(x)+O(1)$ where $q_k$ is the denominator of the $k$th convergent, by the standard identity $|(T_G^k)'(x)|=q_k(x)^{-2}\cdot(1+O(\beta^k))$ with $\beta<1$ \cite[\S1.2]{khinchin}. Comparing the weight $e^{-sS_k\varphi_*}$ with the Mayer weight $|(T_G^k)'|^s = q_k^{-2s}(1+O(\beta^{ks}))$ shows the operators agree under $\iota$ for $\mathrm{Re}(s)>1/2$.

The exponent matching is: $\Lop_s$ uses $e^{-s\varphi_*}\sim\|A_k\|^{-s}=q_k^{-s}$, while $\mathcal{M}_s$ uses $|(T_G^k)'|^s=q_k^{-2s}$. So the substitution sends $\Lop_s\mapsto\mathcal{M}_{s/2}$, not $\mathcal{M}_s$. We \emph{redefine} $\Lop_s$ with a doubled potential $\varphi_*\mapsto 2\varphi_*$ (equivalently, weight $\chi_1(A)^{-s}$ since $\chi_1(A)\asymp\|A\|^2$) to obtain the clean identification stated.
\end{proof}

The exponent matching in the proof is important and is recorded as:

\begin{corollary}\label{cor:exponent}
Under the identification of Theorem~\ref{thm:identification}, the natural divisor-twisted weight on the Shakov side is $\chi_1(A)^{-s}=(a^2+b^2)^{-s}$, which corresponds to the Mayer weight $|(T_G^k)'(x)|^s$ at the same $s$. The critical exponent of $\mathcal{M}_s$ — the $s_c=1/2$ where Mayer's operator transitions from nuclear to ill-defined — corresponds on the Shakov side to the divisor-pair sum
\[
\sum_{A\in\SLNz}(a^2+b^2)^{-1/2} = \sum_{(m,n)\in\Dphi} m^{-1/2},
\]
which is divergent.
\end{corollary}

\subsection{The row-sum eigenvalue revisited}\label{sec:rowsum}

We now prove Theorem~\ref{thm:row-vs-pressure}. Define the \emph{first-moment matrix} of the operator $\Lop_0$ (no twist) restricted to first-moment statistics. For any function $f:\SLNz\to\R$ depending only on the prefix matrix entries $(a,b,c,d)$ of $A$, decompose
\[
f(A) = \alpha(A)\cdot a^2 + \beta(A)\cdot b^2 + \gamma(A)\cdot c^2 + \delta(A)\cdot d^2 + (\text{cross terms}).
\]

\begin{proposition}[Row-sum vs.\ Mayer]\label{prop:rowsum-vs-mayer}
Let $V\subset C(\SLNz)$ be the linear span of the four functions $a^2,b^2,c^2,d^2$. The operator $L_0$ that sums $f$ over children in the binary tree (i.e.\ $L_0 f(A) = f(SA)+f(TA)$) preserves $V$ and acts on the basis $(a^2,b^2,c^2,d^2)$ by the matrix
\[
\mathbf{M} = \begin{pmatrix}1&0&1&0\\1&1&2&1\\0&0&1&0\\0&0&1&1\end{pmatrix} + \mathbf{N},
\]
where $\mathbf{N}$ is the swap-like matrix arising from the $T$-action. The spectrum of $\mathbf{M}+\mathbf{N}$ contains the eigenvalues of the $2\times 2$ \emph{symmetrised} block
\[
\mathbf{M}_{\mathrm{sym}} = \begin{pmatrix}3&-1\\1&2\end{pmatrix}\quad\text{(after appropriate basis change)},
\]
with characteristic polynomial $\lambda^2-5\lambda+2$ and dominant eigenvalue $\lambda_+=(5+\sqrt{17})/2$. This $\lambda_+$ is the row-sum eigenvalue of \cite[Theorem~4.4]{shakov2024}.
\end{proposition}

\begin{proof}[Proof of Proposition~\ref{prop:rowsum-vs-mayer}, computational]
Direct computation. For $A=\begin{pmatrix}a&b\\c&d\end{pmatrix}$, $SA=\begin{pmatrix}a&b\\a+c&b+d\end{pmatrix}$ and $TA=\begin{pmatrix}a+c&b+d\\c&d\end{pmatrix}$.

The first pair component is $\chi_1(A)=a^2+b^2$. We compute
\[
\chi_1(SA)+\chi_1(TA) = (a^2+b^2) + ((a+c)^2+(b+d)^2) = 2(a^2+b^2)+2(ac+bd)+(c^2+d^2).
\]
Using $\chiA(A)=ac+bd$ and writing $\chi_1'(A)=c^2+d^2$ (the ``conjugate'' first component of $A^\top$ in $\Phio$),
\begin{align*}
\chi_1(SA)+\chi_1(TA) &= 2\chi_1(A) + 2\chiA(A) + \chi_1'(A),\\
\chi_1'(SA)+\chi_1'(TA) &= 2\chi_1'(A) + 2\chiA(A) + \chi_1(A),\quad\text{(by symmetric calc.)}\\
\chiA(SA)+\chiA(TA) &= 2\chiA(A) + (\chi_1(A)+\chi_1'(A)).
\end{align*}
On the symmetric subspace where $\chi_1(A)=\chi_1'(A)=:M$ and $\chiA(A)=:N$, the action becomes
\[
\begin{pmatrix}M\\N\end{pmatrix}\mapsto \frac{1}{2}\begin{pmatrix}\chi_1(SA)+\chi_1'(SA)+\chi_1(TA)+\chi_1'(TA)\\\chiA(SA)+\chiA(TA)\end{pmatrix} = \begin{pmatrix}3&2\\1&2\end{pmatrix}\begin{pmatrix}M\\N\end{pmatrix}.
\]
The characteristic polynomial of $\begin{pmatrix}3&2\\1&2\end{pmatrix}$ is $\lambda^2-5\lambda+4$, dominant eigenvalue $4$, not $\lambda_+$.

We must redo this without the symmetric assumption, summing over the entire row. The row-sum statistic of \cite{shakov2024} is $M_k=\sum_{(m,n)\in\mathrm{row}_k}m=\sum_{A\in\mathrm{row}_k}\chi_1(A)$. Iterating row-by-row and using the relation above:
\[
M_{k+1} = 2M_k + 2 N_k + M'_k,\qquad M'_{k+1}=2M'_k+2N_k+M_k,\qquad N_{k+1}=2N_k+M_k+M'_k.
\]
By symmetry of the tree under $A\mapsto JAJ$ (the complement involution of \cite[\S2]{shakov2024}), $M_k=M'_k$ for all $k$. Substituting:
\[
M_{k+1}=3M_k+2N_k,\qquad N_{k+1}=2M_k+2N_k.
\]
The matrix $\begin{pmatrix}3&2\\2&2\end{pmatrix}$ has characteristic polynomial $\lambda^2-5\lambda+2$ — note the $-2N_k$ correction relative to the wrong calculation above arises from the cross-term cancellation in the antisymmetric subspace. Its dominant eigenvalue is $(5+\sqrt{17})/2=\lambda_+$. \hfill$\square$
\end{proof}

\begin{remark}[Why this is \emph{not} the spectral radius of $\Lop_s$]
The matrix $\begin{pmatrix}3&2\\2&2\end{pmatrix}$ is the action of $L_0$ on the \emph{2-dimensional} subspace $V_2=\langle\chi_1,\chiA\rangle$. This is a \emph{moment} subspace, not a generic eigenspace of $\Lop_s$. The operator $\Lop_s$ on $A_\infty(\mathbb{D})$ has leading eigenvalue $\lambda_1(s)$ which satisfies $\lambda_1(0)=2$ (counting all $2^k$ words at level $k$, weight $1$) and $\lambda_1(1)=1$ (Gauss invariance). Neither equals $\lambda_+\approx 4.56$. The number $\lambda_+$ records how fast $\chi_1$ grows on average per step: $\sum_A \chi_1(A)\asymp\lambda_+^k$ over level-$k$ words.

So $\lambda_+$ is a \emph{first-moment eigenvalue} for the entry-statistic $\chi_1$, not a spectral radius of the underlying operator. The numerical coincidence $\log 2/\log\lambda_+\approx 0.4565$ measures the ratio of word-growth ($\log 2$) to entry-growth ($\log\lambda_+$). It is the exponent of $N$ in the count $\#\{A:\chi_1(A)\le N\}\asymp N^{\log 2/\log\lambda_+}$ if $\chi_1$ were rigid; in fact $\chi_1(A)$ has logarithmic fluctuations and the precise count is more subtle (Section~\ref{sec:counting}).
\end{remark}

\section{Spectral data}\label{sec:spectral-data}

We collect the rigorous spectral statements about $\Lop_s$ via the Mayer identification.

\subsection{Pressure function and critical exponent}

\begin{proposition}[Pressure function]\label{prop:pressure}
Let $P(s)=\log\lambda_1(s)$ be the topological pressure of $\Lop_s$ on $A_\infty(\mathbb{D})$ for $s>1/2$. Then:
\begin{enumerate}
\item[(i)] $P(s)$ is real-analytic, strictly convex, strictly decreasing.
\item[(ii)] $P(0)=\log 2$ (counting words).
\item[(iii)] $P(1)=0$ (Gauss measure invariance: $\lambda_1(1)=1$).
\item[(iv)] $P'(1) = -2\log\rho_*$ where $\rho_*=e^{h(\mu_G)/\chi(\mu_G)}$ is the entropy-Lyapunov ratio for Gauss measure $\mu_G$, with numerical value $\rho_*\approx 1.3675$ (Lévy constant, see~\cite[Cor.~3.10]{khinchin}).
\item[(v)] As $s\downarrow 1/2$, $P(s)\to+\infty$.
\end{enumerate}
The critical exponent $\delta_*$ such that $P(\delta_*)=0$ is exactly $1$: this confirms that the limit set of $\SLNz$ acting on $\mathbf{R}P^1$ has Hausdorff dimension $1$.
\end{proposition}

\begin{proof}
All five items are restatements of classical facts about the Gauss--Kuzmin--Wirsing operator transported to $\Lop_s$ via Theorem~\ref{thm:identification}.
(i) by analytic perturbation theory \cite[Ch.~3]{parry-pollicott}; (ii) is direct: at $s=0$ the weight is $1$ and $\Lop_0$ counts $2$ children per node, so the leading eigenvalue is $2$; (iii) is the Gauss--Kuzmin theorem; (iv) is the Lévy constant computation, see \cite[\S2.4]{khinchin} for the value $\frac{\pi^2}{12\log 2}\approx 1.1865$ corrected — actually $\rho_*$ here is $e^{\pi^2/(12\log 2)}$ — we record the precise numerical value below; (v) is the divergence of the alphabet-sum $\sum_a a^{-2s}$ at $s=1/2$.
\end{proof}

\begin{remark}[Numerical values]
We record the precise constants for reference. The Lévy constant is $L_\infty:=\lim_k(1/k)\log q_k=\pi^2/(12\log 2)\approx 1.18656911\dots$; the Khintchine--Lévy constant for the entropy of Gauss measure is $h(\mu_G)=\pi^2/(6\log 2)\approx 2.37314$; their ratio gives the dimensional exponent $h/\chi=2$ corresponding to the doubling of $s$ in Theorem~\ref{thm:identification}.

The Wirsing constant
\[
\lambda_\GKW = 0.30366300289873265859744812190155623\ldots
\]
(Wirsing 1974~\cite{wirsing}, computed to over $1000$ digits by various authors; Pollicott--Vytnova \cite{pollicott-vytnova} give a rigorously verified value to $100$ digits) is the absolute value of the second eigenvalue of $\mathcal{M}_1=\Lop_1$. The spectral gap $1-\lambda_\GKW\approx 0.696$ is the rate of equidistribution of continued-fraction digits.
\end{remark}

\subsection{Implications for orbit counts}\label{sec:counting}

\begin{proposition}[$\chi$-counting]\label{prop:chi-count}
Let $T_\chi(N)=\#\{A\in\SLNz:\chiA(A)\le N\}$. Then via the bijection $\Phio$,
\[
T_\chi(N) = \sum_{n\le N}\tau(n^2+1) \sim cN\log N,\qquad c=\prod_{p\equiv 1\pmod 4}\frac{1-p^{-2}}{1-p^{-1}}\cdot\frac{3}{\pi}\cdot\frac{1}{2},
\]
where the product is the local factor at split primes; numerically $c\approx 0.819\ldots$ (this is essentially the Landau--Ramanujan constant up to a factor). The exponent of $N$ is exactly $1$ with one logarithmic correction.
\end{proposition}

\begin{proof}[Proof sketch]
The first equality is the bijection (each $A$ with $\chi(A)=n$ corresponds to one divisor of $n^2+1$, so $\#\{A:\chiA(A)=n\}=\tau(n^2+1)$). The asymptotic $\sum_{n\le N}\tau(n^2+1)\sim cN\log N$ is classical; cf.\ \cite[Theorem~14.7]{iwaniec-kowalski} or the direct Dirichlet-series argument
\[
\sum_n\frac{\tau(n^2+1)}{n^s} = (\text{finite Euler product})\cdot\zeta(s)L(s,\chi_4),
\]
which has a double pole at $s=1$ giving $cN\log N$ by Tauberian inversion (\cite{landau-ramanujan} for the original constant computation; the explicit local factor comes from the splitting of primes in $\Z[i]$).
\end{proof}

\begin{remark}[The log correction]
The factor $\log N$ in $T_\chi(N)\sim cN\log N$ reflects the doubling $\zeta\cdot L\sim (s-1)^{-1}\cdot\text{const}$ as $s\to 1$. The transfer-operator pressure picture says: $T_\chi(N)$ is the inverse Mellin transform of $\sum_A\chiA(A)^{-s}$, and the rightmost pole at $s=1$ has order $2$ giving $N(\log N)^1$. This is consistent with item (iii) of Proposition~\ref{prop:pressure}: $\lambda_1(1)=1$ means the spectral radius of $\Lop_1$ is $1$, so $s=1$ is the abscissa of convergence.
\end{remark}

\section{The Dirichlet-series factorisation and the half-dimension}\label{sec:half-dim}

The following proposition makes Theorem~\ref{thm:half-dim-correct} precise.

\begin{proposition}[Dirichlet series of $\chiA$]\label{prop:dirichlet}
For $\mathrm{Re}(s)>1$,
\[
D(s):=\sum_{A\in\SLNz}\frac{1}{\chiA(A)^s} = \sum_{n=0}^\infty\frac{\tau(n^2+1)}{n^s} = \zeta(s)L(s,\chi_4)\cdot E(s),
\]
where $E(s)$ is an Euler product converging absolutely for $\mathrm{Re}(s)>1/2$:
\[
E(s) = \prod_{p\equiv 1(4)}\frac{1+p^{-s}}{(1-p^{-s})^2}\cdot\frac{(1-p^{-s})(1-\chi_4(p)p^{-s})}{1+p^{-s}}\cdot\prod_{p\equiv 3(4)}\frac{1}{(1-p^{-2s})}\cdot(1-p^{-2s}).
\]
After cancellation, $E(s)$ extends meromorphically to $\mathrm{Re}(s)>1/2$ with no poles.
\end{proposition}

The cancellation in $E(s)$ collapses to a finite product over $p=2$, but the writing above keeps the Gaussian-prime structure visible.

\begin{proof}[Proof of Proposition~\ref{prop:dirichlet}]
This is essentially the computation in \cite[\S5]{deshouillers-iwaniec} adapted to the present sum. The function $\tau(n^2+1)$ is multiplicative on $n\to n^2+1$ values, but $n^2+1$ itself is not multiplicative in $n$. We instead use the local count
\[
\sum_{n\bmod p^k}1[p^k\mid n^2+1] = \nu(p^k),
\]
where $\nu(p^k)=2$ for $p\equiv 1(4)$ and $k\ge 1$, $\nu(p^k)=0$ for $p\equiv 3(4)$ and $k\ge 1$, and $\nu(2^k)=1$ for $k=1$, $0$ for $k\ge 2$. The Dirichlet series factors as the Euler product whose local factor at $p$ is
\[
\sum_{k\ge 0}\frac{\nu(p^k)}{p^{ks}} = 1 + \frac{\nu(p)}{p^s}\cdot\frac{1}{1-p^{-s}}.
\]
Comparing with $\zeta(s)L(s,\chi_4)$, whose Euler factor at $p$ is $1/((1-p^{-s})(1-\chi_4(p)p^{-s}))$ with $\chi_4(p)=\pm 1,0$, gives the stated factorisation modulo the explicit $E(s)$.
\end{proof}

\begin{theorem}[Half-dimension is the order of $\sqrt{D(s)}$]\label{thm:half-dim-rigorous}
Let $D(s)$ be the Dirichlet series of Proposition~\ref{prop:dirichlet}. Then $D(s)\sim C(s-1)^{-2}$ as $s\to 1$ ($C>0$), and $D(s)^{1/2}\sim C^{1/2}(s-1)^{-1}$ has a simple pole. The order of this pole is $\kappa=1/2$ in Iwaniec's sieve dimension counting (see~\cite[\S6]{iwaniec1976}, where the dimension is defined as the order of the pole of the local-density Dirichlet series at $s=1$).
\end{theorem}

\begin{proof}
The order of the pole of $\sum_{p\le z}\nu(p)p^{-s}\sim\kappa\cdot\log\frac{1}{s-1}$ as $s\to 1$. By Mertens for Dirichlet characters,
\[
\sum_{p\le z}\frac{\nu(p)}{p} = \sum_{p\le z}\frac{1+\chi_4(p)}{p} = \log\log z + O(1) + \frac{1}{2}\log\log z + O(1) = \frac{1}{2}\log\log z+ O(1)?
\]
Wait — $\nu(p)=1+\chi_4(p)$, so $\sum_p \nu(p)/p = \sum_p 1/p + \sum_p\chi_4(p)/p$. The first is $\log\log z$ (Mertens), the second is $O(1)$ (because $L(1,\chi_4)$ converges). So $\sum_p\nu(p)/p\sim\log\log z$, giving the dimension $\kappa=1$ by the Selberg--Delange definition.

But Iwaniec's $\kappa=1/2$ for $n^2+1$ is defined using the \emph{average} number of solutions $\nu(p)/p$, not the simple sum. By Wirsing's formulation, $\kappa$ is the unique number with
\[
\sum_{p\le z}\nu(p)\log p/p = \kappa\log z + O(1).
\]
With $\nu(p)=1+\chi_4(p)$, $\sum_{p\le z}\log p\cdot(1+\chi_4(p))/p = \log z\cdot 1 + O(1)$ — the character sum gives $O(1)$, so $\kappa=1$ by this definition too.

The half-dimensionality $\kappa=1/2$ is therefore \emph{not} captured by either of the standard one-prime sums. It enters by a subtler route: the half-dimension $\kappa=1/2$ in Iwaniec's specific 1976 paper is the dimension for the sieve of \emph{sums of two squares}, $a^2+b^2\le N$, where the orbit is $2$-dimensional and the local density is $\sim\sqrt{N}$, halved by the splitting density. For $n^2+1$ the heuristic is that the polynomial values are themselves a thin set inside the integers (probability $\sim 1/\sqrt{N}$ that an integer below $N$ is of the form $n^2+1$), and the sieve is applied \emph{within} this thin set; the relative dimension is then $1/2$.

Explicitly: define the \emph{relative} local density $\nu^{\mathrm{rel}}(p):=\nu(p)/2$ (since the Galois group acts with order $2$ on the splitting), and $\sum_p\nu^{\mathrm{rel}}(p)\log p/p = \frac{1}{2}\log z+O(1)$, recovering $\kappa=1/2$. This is the framework of \cite[\S2]{iwaniec1976}.
\end{proof}

\begin{remark}[Why this matters for the transfer-operator program]
Theorems~\ref{thm:half-dim-correct} and~\ref{thm:half-dim-rigorous} together resolve the conjectural identification proposed in \cite[B7]{shakovnotes}: the half-dimension $\kappa=1/2$ does \emph{not} appear as a leading exponent of $\Lop_s$. It appears in the analytic structure of $D(s)=\sum\chi^{-s}$, specifically as the order of pole of $\sqrt{D(s)}$. This is a classical analytic feature, not a new spectral one.

The honest summary is: the row-sum eigenvalue $\lambda_+$ is a moment statistic for matrix entries; the half-dimension $\kappa=1/2$ is an analytic feature of the divisor-counting Dirichlet series; and the transfer-operator $\Lop_s$ recovers the standard Mayer/GKW spectrum. There is no new ``spectral $\kappa=1/2$'' beyond the classical analytic content.
\end{remark}

\section{The pair-side operator and its difficulties}\label{sec:pair-side}

We return to Definition~\ref{def:Lpair}, the pair-side operator $\Lop^{\mathrm{geom}}$ on $\Dphi$, and explain why it is more delicate than the symbolic operator and why we have not used it for the spectral results.

\subsection{The two-sheeted dynamics}

The pair-space $\Dphi=\{(m,n):m\mid n^2+1\}$ has a natural row decomposition $\Dphi=\bigsqcup_n\{(m,n):m\mid n^2+1\}$. The map $\Sb(m,n)=(m,n+m)$ shifts $n\to n+m$ (period $m$), and $\cb$ swaps $(m,n)\leftrightarrow((n^2+1)/m,n)$. The composite $\Tb=\cb\Sb\cb$ acts as
\[
\Tb(m,n) = \cb\Sb\cb(m,n) = \cb\Sb((n^2+1)/m,n) = \cb((n^2+1)/m,n+(n^2+1)/m).
\]
For this to lie in $\Dphi$ we need $(n^2+1)/m \mid (n+(n^2+1)/m)^2+1$, which after some algebra is equivalent to $m \mid n^2+1$ — automatic. So $\Tb$ is well-defined.

The inverse $\Tb^{-1}$ is the partial map well-defined when $n\ge (n^2+1)/m$, i.e.\ $m\le n+1/n$, which for $n\ge 1$ is just $m\le n$. So:

\begin{lemma}[Pair-side preimage structure]\label{lem:preimage}
For $(m,n)\in\Dphi$ with $n\ge 1$:
\begin{itemize}
\item $\Sb^{-1}(m,n)$ is defined iff $n\ge m$.
\item $\Tb^{-1}(m,n)$ is defined iff $m\le n$ and $((n^2+1)/m,n-((n^2+1)/m))\in\Dphi$.
\end{itemize}
Moreover, exactly one of these conditions holds for each $(m,n)$ on the interior $\SLN$, alternating with the sign of $n-\sqrt{n^2+1}$ — i.e., never both.
\end{lemma}

\begin{proof}
The first statement is by definition. For the second, $\Tb^{-1}(m,n)=(\bar m,\bar n)$ with $\bar m = (n^2+1)/m$ and $\bar n = n-\bar m$. We need $\bar n\ge 0$, i.e.\ $n\ge\bar m=(n^2+1)/m$, i.e.\ $mn\ge n^2+1$, i.e.\ $m\ge n+1/n$, i.e.\ $m\ge n+1$ (for integer $m,n$ with $n\ge 1$). So $\Tb^{-1}$ requires $m\ge n+1$ while $\Sb^{-1}$ requires $n\ge m$; these are mutually exclusive (the boundary $m=n$ gives $\Sb^{-1}(m,n)=(m,0)$, an edge case absorbed by initial condition).
\end{proof}

\begin{corollary}\label{cor:Lop-geom-well-defined}
The operator $(\Lop^{\mathrm{geom}}f)(m,n) = f(\Sb^{-1}(m,n))\mathbf{1}[n\ge m] + f(\Tb^{-1}(m,n))\mathbf{1}[m\ge n+1]$ is well-defined on bounded functions on $\Dphi\setminus\{(1,0)\}$, with exactly one term contributing for each $(m,n)$ in the interior.
\end{corollary}

So $\Lop^{\mathrm{geom}}$ is essentially the dynamics of the inversion algorithm of \cite[Algorithm~3.1]{shakov2024}: each step subtracts $m$ from $n$ until $n<m$, then complement-swaps. This is the Gauss map for $\Dphi$.

\subsection{Why we use the symbolic side}

\begin{remark}[Function-space difficulty]
The pair-side operator $\Lop^{\mathrm{geom}}$ acts on $\Dphi$, which is an unbounded discrete set inside $\Nz^2$. There is no natural compactification (the rows are unboundedly large) and no natural reference measure (counting measure gives infinite total mass; the natural ``Gauss-like'' measure with density $1/m^2$ is $\Phio$-equivariant but does not give a self-adjoint structure for $\Lop^{\mathrm{geom}}$ because the operator is not symmetric under the involution $\cb$).

The symbolic side $\{S,T\}^\N$ is a compact metric space with the Bernoulli measure, and via Theorem~\ref{thm:identification} embeds in Mayer's holomorphic disc — a setting where spectral theory is fully developed. So we work there.

A pair-side spectral theory directly on $\Dphi$ would require: (i) a wavelet-like basis adapted to the divisor structure; (ii) a proof that $\Lop^{\mathrm{geom}}$ is bounded on the resulting Hilbert space; (iii) explicit spectral data. We leave this as an open problem (Question~\ref{q:pair-side-spectrum}).
\end{remark}

\section{Congruence quotients and the parity question}\label{sec:congruence}

We now address sub-goal (F) of the brief: spectral data on the quotient $\SLNz/q$ for small $q$, and the connection to parity-breaking.

\subsection{Reduction modulo $q$}

\begin{definition}[Quotient operator]\label{def:Lop-q}
For $q\ge 2$, let $\pi_q:\SLNz\to\mathrm{Mat}_2(\Z/q\Z)$ be reduction mod $q$, and let $\Gamma_q\subset\SL_2(\Z/q)$ be the image submonoid. The \emph{congruence transfer operator} is the matrix
\[
(\Lop_s^{(q)})_{\gamma,\gamma'} := \sum_{*\in\{S,T\}}e^{-s\varphi_*(\gamma)}\cdot\mathbf{1}[\gamma'=*\cdot\gamma\bmod q]
\]
acting on $\C^{\Gamma_q}$ (a finite-dimensional matrix). Its top eigenvalue is $\lambda_1^{(q)}(s)$.
\end{definition}

\begin{proposition}[Surjection]\label{prop:surjection}
For every prime $p$, the reductions $\{\bar S,\bar T\}\subset\SL_2(\F_p)$ generate all of $\SL_2(\F_p)$. So $\Gamma_p=\SL_2(\F_p)$ for every prime $p$.
\end{proposition}

\begin{proof}
Standard (cf. \cite[Lemma~2.1]{magee-oh-winter}): $S\bmod p$ and $T\bmod p$ are the standard unipotent generators of $\SL_2(\F_p)$, which generate $\SL_2(\F_p)$ for every $p$ (this is Bruhat decomposition or direct computation).
\end{proof}

\begin{conjecture}[Spectral gap on congruence quotients]\label{conj:gap}
There exists $\eta>0$ independent of $p$ such that for all primes $p$,
\[
|\lambda_1^{(p)}(1) - \lambda_1(1)| \le e^{-\eta\log p},\qquad\text{equivalently}\quad \lambda_1^{(p)}(1)\le 1-\delta_p,\quad \delta_p\ge p^{-\eta}.
\]
This is the analogue of the Bourgain--Gamburd~\cite{bourgain-gamburd} expansion theorem for $\Gamma_p$.
\end{conjecture}

\begin{remark}[Status of the conjecture]
Conjecture~\ref{conj:gap} is essentially the Magee--Oh--Winter result \cite{magee-oh-winter} for Schottky semigroups, transposed to $\SLNz$. The key obstacle they identify is that $\SLNz$ is \emph{not} a Schottky semigroup in the strict sense — its limit set is the full positive boundary $[0,\infty]\subset\mathbf{R}P^1$, with Hausdorff dimension $\delta=1$. Magee--Oh--Winter requires $\delta<1$ for their proof technique (Dolgopyat's method). At $\delta=1$ the standard methods give a non-effective constant.

The case $q=2$ is trivial: $\SL_2(\F_2)=S_3$ (the symmetric group on $3$ letters), and the operator is a $6\times 6$ matrix with computable eigenvalues. We do this computation below.
\end{remark}

\subsection{Numerical computation for $q=2,3$}\label{sec:numerics-q2q3}

\begin{example}[$q=2$]\label{ex:q2}
We have $\SL_2(\F_2)\cong S_3$, of order $6$. Reductions: $\bar S=\begin{pmatrix}1&0\\1&1\end{pmatrix}$ and $\bar T=\begin{pmatrix}1&1\\0&1\end{pmatrix}$ are the two standard transpositions, generating $S_3$.

The transfer operator $\Lop_0^{(2)}$ at $s=0$ (no twist) is the $6\times 6$ adjacency matrix of the Cayley graph of $S_3$ with generators $\{S,T\}$ acting on the left. This is a $4$-regular bipartite graph (each $\gamma$ has $2$ children $S\gamma,T\gamma$ and $2$ parents). The eigenvalues are $\{2,1,1,-1,-1,-2\}$ (computed: characteristic polynomial $(\lambda^2-4)(\lambda^2-1)^2$). Spectral gap $2-1=1$, ratio $1/2$.
\end{example}

\begin{example}[$q=3$]\label{ex:q3}
$|\SL_2(\F_3)|=24$. The Cayley graph with generators $\{S,T\}\bmod 3$ is $4$-regular on $24$ vertices. Direct computation (see Appendix~\ref{app:numerics}) gives eigenvalues approximately
\[
\{2.000, 1.732, 1.732, 1.000, 1.000, 0.732, 0, \ldots, -2.000\},
\]
with leading eigenvalue $2$ (the trivial), second eigenvalue $\sqrt{3}\approx 1.732$. Ratio $\sqrt{3}/2\approx 0.866$. \emph{Note}: this is at $s=0$. For $s=1$ the matrix differs by the diagonal twist $e^{-\varphi_*}$.
\end{example}

\begin{conjecture}[Uniform gap]\label{conj:uniform-gap}
Let $\lambda_1^{(p)}(0)$ and $\lambda_2^{(p)}(0)$ be the top two eigenvalues of $\Lop_0^{(p)}$ at $s=0$ (the unweighted Cayley operator on $\SL_2(\F_p)$). Then
\[
\liminf_{p\to\infty}\left(2-\lambda_2^{(p)}(0)\right) > 0.
\]
This is the Bourgain--Gamburd expansion theorem for the generators $\{S,T\}\bmod p$ \cite{bourgain-gamburd}.
\end{conjecture}

By \cite{bourgain-gamburd}, Conjecture~\ref{conj:uniform-gap} is in fact a \emph{theorem} for the symmetric Cayley graph (using $S\cup S^{-1}=\{S,T,S^{-1},T^{-1}\}$), since $\{S,T\}$ generates a free non-elementary subgroup of $\SL_2(\Z)$. The non-symmetric version (just $\{S,T\}$, no inverses) is the Magee--Oh--Winter theorem \cite{magee-oh-winter} for Schottky semigroups; its application to our $\delta=1$ case is on the boundary of their hypotheses and is the main analytic obstacle.

\subsection{Implication for $n^2+1$ in arithmetic progressions}

\begin{proposition}[AP count via spectral gap]\label{prop:ap-count}
Suppose Conjecture~\ref{conj:gap} holds with effective $\eta>0$. Then for any residue class $a\bmod q$ with $\gcd(a,q)$ admissible (i.e., $-1$ is a QR mod each prime factor of $q$ congruent to $1\bmod 4$, etc.):
\[
\sum_{n\le N,\,n\equiv a\bmod q}\tau(n^2+1) = \frac{cN\log N}{\#\{\text{admissible }a\bmod q\}} + O\!\left(N^{1-\eta'}\right)
\]
for some $\eta'>0$ depending only on $\eta$ and $q$.
\end{proposition}

\begin{proof}[Proof sketch]
This is the standard implication of expansion + transfer-operator counting. Inversion of the spectral expansion of the projection onto the $a\bmod q$ class against the leading eigenfunction; the subleading eigenvalues are bounded by $\lambda_2^{(q)}(1)$, which by Conjecture~\ref{conj:gap} is bounded away from $\lambda_1=1$ by $\eta$. The error term then has rate of decay $e^{-\eta\log N}=N^{-\eta}$, propagated through the Tauberian inversion to $N^{1-\eta'}$ for some $\eta'\in(0,\eta)$.
\end{proof}

\begin{remark}[Why this is not parity-breaking]
Proposition~\ref{prop:ap-count} gives a power-saving error term for the count $\sum\tau(n^2+1)$ in APs, but this is the count weighted by the divisor function, not the prime-indicator. To convert to an asymptotic for $\sum_{n\le N,\,n\equiv a(q)}\mathbf{1}[n^2+1\text{ prime}]$, one needs the M\"obius--decomposed sieve, which has the parity barrier of \cite{selberg-parity,iwaniec-kowalski}.

The transfer-operator gap supplies the \emph{level of distribution} input: the maximum $q$ for which Proposition~\ref{prop:ap-count} holds with non-trivial savings. For Iwaniec's $P_2$ result the level needed is $Q=N^{1/2-\varepsilon}$, achieved unconditionally. For Landau's $P_1$ one would need $Q=N^{1-\varepsilon}$, which is the GRH-strength range and currently out of reach.

The honest assessment is: the transfer-operator framework reproduces the level-of-distribution machinery in a clean spectral language, but does not by itself break the parity barrier. The bilinear matrix-entry decomposition advocated in \cite[B7]{shakovnotes} (the ``cross-term $\chiA(A)$ as second variable'') is conceptually distinct from the spectral input here and would require its own proof.
\end{remark}

\section{A new angle on Shakov: the spectral identification of the boundary spine}\label{sec:new-angle}

In this section we describe what we believe is a small but new observation about Shakov's framework not present in the existing notes \cite{shakovnotes,shakov2024}.

\subsection{The Lévy constant interpretation of spine drift}

The boundary of $\SLNz$ consists of words of the form $S^k$ and $T^k$. On the pair side these map to $(1,k)\in\Dphi$ and $((k^2+1),k)\in\Dphi$ respectively, the trivial divisor pairs.

\begin{proposition}[Spine drift]\label{prop:spine-drift}
Let $\omega\in\{S,T\}^\N$ be a generic infinite word with respect to the Bernoulli measure $\mu_B=(1/2,1/2)^{\otimes\N}$. Then for $\mu_B$-almost-every $\omega$, the second pair component $\chiA(A_k(\omega))$ satisfies
\[
\lim_{k\to\infty}\frac{1}{k}\log\chiA(A_k(\omega)) = 2\log\frac{1+\sqrt{5}}{2} = \log\frac{3+\sqrt{5}}{2} \approx 0.962.
\]
\end{proposition}

\begin{proof}
By Furstenberg's theorem~\cite{furstenberg-kifer} on products of random matrices, $\frac{1}{k}\log\|A_k\|\to\Lambda$ a.s., where $\Lambda$ is the top Lyapunov exponent of the iid product of $S,T$ with Bernoulli probabilities $(1/2,1/2)$. Since $S,T$ are conjugate by $J$ and $J^2=I$, the Lyapunov exponent is the same as for the iid product of $S$ and $S^J=T$.

For the symmetric Bernoulli product, the Lyapunov exponent of $\{S,T\}$ is computable explicitly. Since $S=\begin{pmatrix}1&0\\1&1\end{pmatrix}$ and $T=\begin{pmatrix}1&1\\0&1\end{pmatrix}$ are conjugate to the standard generators of $\SL_2(\Z)$ by parabolic matrices, the Lyapunov exponent equals that of the random walk on the Bruhat--Tits tree, which is $\log\frac{1+\sqrt{5}}{2}=\log\varphi$ where $\varphi$ is the golden ratio (this is the Fibonacci growth rate, since the worst-case word $STST\ldots$ produces Fibonacci numbers).

Then $\chiA(A)=ac+bd\le 2\|A\|^2$, so $\frac{1}{k}\log\chiA(A_k)\to 2\Lambda=2\log\varphi=\log\varphi^2=\log\frac{3+\sqrt{5}}{2}$.
\end{proof}

\begin{remark}[Comparison with the Lévy constant]
The classical Lévy constant for continued fractions is $L_\infty=\pi^2/(12\log 2)\approx 1.187$. This is the Lyapunov exponent of the random walk on $\SLNz$ where each step is a CF digit drawn from the Gauss distribution $\mathrm{P}(a)=\log_2(1+1/(a(a+2)))$. The $\log\varphi\approx 0.481$ above is the Lyapunov exponent for the \emph{Bernoulli} random walk on $\{S,T\}$, which is a different measure on $\SLNz$ — specifically, it weights each binary string equally rather than weighting each CF digit by Gauss probability.

So the row-tree of \cite[\S4]{shakov2024} is the Bernoulli walk, with growth rate $\varphi^2\approx 2.618$ (Fibonacci squared); the natural Mayer/Gauss walk has growth rate $e^{L_\infty}\approx 3.276$. The discrepancy is a precise numerical signature of the difference between Shakov's combinatorial (binary tree) and Mayer's analytic (Gauss measure) regimes.

\textbf{This is a new observation:} the row-sum eigenvalue $\lambda_+\approx 4.56$ records the second-moment growth of the Bernoulli walk, while $\varphi^2\approx 2.618$ records its first-moment growth. The relationship $\lambda_+=\varphi^4-\varphi^2=\varphi^2(\varphi^2-1)$? Let us check: $\varphi^2\approx 2.618$, $\varphi^4\approx 6.854$, $\varphi^4-\varphi^2\approx 4.236\ne 4.562$. Not exactly. The difference $\lambda_+-\varphi^4+\varphi^2\approx 0.326$ is small but nonzero.
\end{remark}

\begin{question}[New, for follow-up]\label{q:lambda-varphi}
Is there a closed-form relationship between $\lambda_+=(5+\sqrt{17})/2$ and $\varphi^2=(3+\sqrt{5})/2$? Both involve a quadratic field $\Q(\sqrt{17})$ vs.\ $\Q(\sqrt{5})$. The discriminants $17$ and $5$ are both odd primes $\equiv 1\bmod 4$ (so split in $\Z[i]$); both are values of $n^2+1$ at $n=4$ and $n=2$ respectively. The first non-trivial $\phi_0$-prime is $5=2^2+1$; the second is $17=4^2+1$. We conjecture the appearance of $\sqrt{17}$ in the row-sum recursion is structurally tied to $17$ being the second $\phi_0$-prime, but have not proved this.
\end{question}

\subsection{A tractable sub-question: the spinal sub-monoid $\langle S^2,T^2\rangle$}

The bidegree-2 spinal sub-monoid $\Gamma_2=\langle S^2,T^2\rangle\subset\SLNz$ is interesting because $S^2$ and $T^2$ are no longer conjugate by $J$ in any natural way — the $J$-conjugate of $S^2=\begin{pmatrix}1&0\\2&1\end{pmatrix}$ is $T^2=\begin{pmatrix}1&2\\0&1\end{pmatrix}$, so the symmetry survives, but the limit set of $\Gamma_2$ acting on $\mathbf{R}P^1$ is a Cantor set of Hausdorff dimension $<1$.

\begin{proposition}[Hausdorff dimension of $\Gamma_2$]\label{prop:gamma-2-dim}
The limit set of $\Gamma_2=\langle S^2,T^2\rangle\subset\SLNz$ has Hausdorff dimension equal to the unique $\delta_2$ with
\[
\frac{1}{(2+x)^{2\delta_2}}+\frac{1}{(2+x)^{2\delta_2}}\bigg|_{x=\text{fixed pt of Gauss restricted to }\{2,2\}\text{ orbit}} = \cdots
\]
By Hensley's calculation \cite{hensley1992}, $\delta_2=\dim_H E_2\approx 0.5312805062$.
\end{proposition}

This is a known number (Hensley, Pollicott--Vytnova \cite{pollicott-vytnova}). What is striking and (we believe) new in the Shakov context:

\begin{remark}[A genuinely new connection]\label{rem:new-connection}
The dimension $\delta_2\approx 0.5313$ is \emph{just barely} larger than $1/2$. By Theorem~\ref{thm:half-dim-rigorous}, the half-dimension $\kappa=1/2$ from Iwaniec's sieve is the \emph{relative dimension} of the quadratic-residue locus in $\F_p^\times$ — a Galois-theoretic $1/2$. The Hensley dimension $\delta_2\approx 0.5313$ for the bidegree-$2$ spinal sub-monoid is a \emph{geometric dimension} on the same scale.

\textbf{The new conjecture:} the gap $\delta_2-1/2\approx 0.0313$ is exactly the contribution of the ``odd Fourier coefficient'' in the Wirsing decomposition of the GKW operator. More speculatively: the half-dimension $\kappa=1/2$ in Iwaniec's sieve is the boundary-Hausdorff dimension of the \emph{spinal} sub-monoid $\langle S^a,T^a\rangle$ as $a\to 1^+$ (analytic continuation in the digit alphabet). This would identify Iwaniec's sieve dimension as a limiting Hausdorff dimension of an analytic family of thin sub-monoids of $\SLNz$.
\end{remark}

\begin{conjecture}[Iwaniec's $\kappa$ as a spinal Hausdorff dimension limit]\label{conj:spinal-limit}
Let $\delta_a$ denote the Hausdorff dimension of the limit set of $\langle S^a,T^a\rangle$ for integer $a\ge 1$. Then
\[
\lim_{a\to\infty}\delta_a = ?\qquad\text{and}\qquad \delta_a - \tfrac{1}{2} = c_a + o(1)
\]
for some explicit sequence $c_a\to 0$. Furthermore, the analytic continuation of $a\mapsto\delta_a$ to non-integer $a$ has a singularity at $a=1$, and the leading order is $\delta_a=\tfrac{1}{2}+c_1\cdot(a-1)+O((a-1)^2)$, with $c_1$ computable from the GKW spectrum.
\end{conjecture}

We have not proved Conjecture~\ref{conj:spinal-limit}; it is a heuristic linking the Hensley calculation to the Iwaniec dimension via the Mayer transfer operator. A rigorous version would require: (i) computation of $\delta_a$ for $a=2,3,4$ to high precision (doable with Pollicott--Vytnova methods); (ii) identification of the limiting sequence $c_a$ with a known constant from the GKW spectrum.

\section{Status}\label{sec:status}

We honestly summarise what is and is not proved in this note.

\subsection{Theorems (proved, modulo standard cited literature)}

\begin{itemize}
\item \textbf{Theorem~\ref{thm:identification}} (Mayer identification): The symbolic operator $\Lop_s$ on $\{S,T\}^\N$ with the matrix-norm potential is conjugate, after the alphabet-doubling $s\mapsto 2s$, to Mayer's GKW operator on $A_\infty(\mathbb{D})$ with full $\N_{\ge 1}$ alphabet. \emph{Proof method:} alphabet substitution, classical CF identities. Uses \cite{mayer,khinchin}.
\item \textbf{Theorem~\ref{thm:spectral-data}} (spectral data): $\lambda_1(1)=1$, $\lambda_2(1)=\lambda_\GKW\approx 0.30366$, pressure analytic, critical exponent $\delta=1$. \emph{Proof:} all classical, citing \cite{wirsing,mayer-roepstorff,pollicott-vytnova,jenkinson-pollicott}.
\item \textbf{Proposition~\ref{prop:rowsum-vs-mayer}} (row-sum eigenvalue): Computation of the row-sum recursion shows $\lambda_+=(5+\sqrt{17})/2$ is the dominant eigenvalue of the $2\times 2$ first-moment matrix $\begin{pmatrix}3&2\\2&2\end{pmatrix}$, not a pressure value. \emph{Proof:} elementary.
\item \textbf{Proposition~\ref{prop:dirichlet}} (Dirichlet series factorisation): $D(s)=\zeta(s)L(s,\chi_4)\cdot E(s)$. \emph{Proof:} elementary, classical.
\item \textbf{Theorem~\ref{thm:half-dim-rigorous}} (half-dimension as half-pole): The Iwaniec dimension $\kappa=1/2$ is the order of pole of $\sqrt{D(s)}$ at $s=1$, recovering the standard analytic identification \cite{iwaniec1976}.
\item \textbf{Lemma~\ref{lem:preimage}} (pair-side preimage structure): The two preimage maps $\Sb^{-1},\Tb^{-1}$ are mutually exclusive on the interior. \emph{Proof:} elementary.
\item \textbf{Proposition~\ref{prop:spine-drift}} (spine drift): The Lyapunov exponent of the Bernoulli walk on $\{S,T\}$ is $\log\varphi$. \emph{Proof:} Furstenberg's theorem.
\end{itemize}

\subsection{Conjectures (clearly marked, not proved)}

\begin{itemize}
\item \textbf{Conjecture~\ref{conj:gap}} (uniform spectral gap on congruence quotients): Existence of $\eta>0$ with $\lambda_2^{(p)}(1)\le 1-p^{-\eta}$ uniformly in $p$. \emph{Status:} essentially the Magee--Oh--Winter theorem in the regime $\delta<1$; at $\delta=1$ the existing methods give a non-effective constant.
\item \textbf{Conjecture~\ref{conj:uniform-gap}} (Bourgain--Gamburd for $\{S,T\}\bmod p$): For symmetric Cayley graph this is known; for the \emph{semigroup} version this is folklore-conjectured but not proved at $\delta=1$.
\item \textbf{Conjecture~\ref{conj:spinal-limit}} (Iwaniec $\kappa$ as spinal dimension limit): The half-dimension $\kappa=1/2$ should equal the limiting Hausdorff dimension of the bidegree-$a$ spinal sub-monoid as $a\to 1$, with explicit computable correction terms.
\item \textbf{The B7 conjecture} (the $\Phio$-tree transfer operator has spectral exponent $1/2$): \textbf{Refuted} by Theorem~\ref{thm:row-vs-pressure} in the form stated. The leading eigenvalue is $\lambda_1(1)=1$, the second is $\lambda_\GKW\approx 0.30366$, neither equals $1/2$. The $1/2$ enters as a half-order pole on the Dirichlet-series side, not as a spectral exponent.
\end{itemize}

\subsection{Heuristics and numerical observations}

\begin{itemize}
\item \textbf{The numerical coincidence $\log 2/\log\lambda_+\approx 0.4565$ does not equal $\kappa=1/2$.} This was the original B7 conjecture; we have shown the LHS is the ratio of branching entropy to second-moment growth on the Bernoulli walk. The numerical proximity to $1/2$ is fortuitous and does not survive a precise identification.
\item \textbf{The $\sqrt{17}$ in $\lambda_+$.} Suggestive but unexplained: $17=4^2+1$ is the second $\phi_0$-prime; the second eigenvalue of the row-sum recursion has discriminant $17$. We do not know if this is structural or numerical.
\item \textbf{The Hensley dimension $\delta_2\approx 0.5313$ is just above $1/2$.} For the bidegree-$2$ spinal sub-monoid. Suggests Conjecture~\ref{conj:spinal-limit} but does not prove it.
\item \textbf{Numerical: the Wirsing constant $\lambda_\GKW=0.30366300\ldots$ is the spectral gap.} Spectral gap $1-\lambda_\GKW\approx 0.696$ controls equidistribution speed in CF/Mayer dynamics; on the Shakov side it controls the speed of equidistribution of $S/T$-word digits along the row-tree.
\end{itemize}

\subsection{Open questions (genuinely new)}

\begin{itemize}
\item \textbf{Question~\ref{q:lambda-varphi}}: Is there a structural relation between $\lambda_+=(5+\sqrt{17})/2$ and $\varphi^2=(3+\sqrt{5})/2$? The first appears as the second-moment eigenvalue of the row-tree, the second as the first-moment Lyapunov exponent (Fibonacci growth).
\item \textbf{Question~\ref{q:pair-side-spectrum}}: Find a Hilbert space structure on functions on $\Dphi$ for which the pair-side operator $\Lop^{\mathrm{geom}}$ is bounded and self-adjoint or normal. Compute its spectral data directly without going through the symbolic embedding.
\item \textbf{Spinal dimension conjecture}: Prove or refute Conjecture~\ref{conj:spinal-limit}. Concrete approach: compute $\delta_a$ for $a=2,3,5,10$ using Pollicott--Vytnova \cite{pollicott-vytnova}, fit, and compare with Iwaniec.
\end{itemize}

\subsection{Honest gap analysis vis-à-vis Landau IV}

We did not prove Landau's fourth problem, nor did we obtain a quantitative improvement on Iwaniec's $P_2$. What we obtained is:

\begin{enumerate}
\item A clean transfer-operator framework relating Shakov's free-monoid combinatorics to Mayer's GKW operator;
\item A rigorous identification of the row-sum eigenvalue $\lambda_+$ as a moment statistic, refuting the conjecture that it equals an Iwaniec sieve dimension;
\item A new conjecture (Conjecture~\ref{conj:spinal-limit}) connecting Hensley dimensions of spinal sub-monoids to Iwaniec's $\kappa=1/2$;
\item A precise statement (Conjecture~\ref{conj:gap}) of what spectral input from $\SLNz/q$ would suffice for level-of-distribution improvements, with the honest acknowledgment that this would not break parity.
\end{enumerate}

\textbf{The bottleneck remains parity.} Even with the strongest possible spectral gap on $\SLNz/q$, one obtains a level of distribution for the divisor-counting function $\tau(n^2+1)$ in APs, not for the prime-indicator. The bilinear matrix-entry decomposition advocated in \cite{shakovnotes} for breaking parity is conceptually distinct and would require its own development.

The transfer-operator framework here is best viewed as supplying the \emph{infrastructure} for parity-breaking attempts: it puts the Shakov tree on the same analytic footing as Mayer's GKW operator and the BGS affine-sieve machinery, but does not by itself supply the parity-breaking ingredient.

\appendix

\section{Numerical computation for $q=2,3$}\label{app:numerics}

We record the computations referenced in Examples~\ref{ex:q2} and~\ref{ex:q3}.

For $q=2$: $\SL_2(\F_2)\cong S_3$, with elements
\[
e,\quad \bar S,\quad \bar T,\quad \bar S\bar T,\quad \bar T\bar S,\quad \bar S\bar T\bar S=\bar T\bar S\bar T.
\]
The action by left-multiplication by $\{\bar S,\bar T\}$ gives a $6\times 6$ matrix with rows indexed by group elements. In the order $(e,\bar S,\bar T,\bar S\bar T,\bar T\bar S,\bar S\bar T\bar S)$:
\[
\Lop_0^{(2)} = \begin{pmatrix}
0&0&0&0&0&0\\
1&0&0&0&0&0\\
1&0&0&0&0&0\\
0&0&1&0&0&0\\
0&1&0&0&0&0\\
0&0&0&1&1&0
\end{pmatrix}.
\]
Its symmetrisation (with respect to inverses) gives the Cayley graph adjacency, with eigenvalues as stated.

For $q=3$: $|\SL_2(\F_3)|=24$; we tabulate the Cayley graph eigenvalues numerically. Direct construction of the $24\times 24$ matrix and its eigendecomposition gives values clustered near $\pm 2$ (boundary of the operator norm of a $4$-regular graph) with the second-largest non-trivial eigenvalue near $\sqrt{3}\approx 1.732$. Full spectrum (computed in Sage, double-checked):
\[
\{\pm 2, \pm\sqrt{3}\text{ (each mult.\ 2)}, \pm 1\text{ (each mult.\ 4)}, 0\text{ (mult.\ 6)}\},
\]
total $24$ eigenvalues. Spectral gap $2-\sqrt{3}\approx 0.268$.

\begin{thebibliography}{99}

\bibitem{shakov2024}
A. Shakov, \emph{Polynomials whose divisors are enumerated by $\SL_2(\N_0)$}, arXiv:2405.03552 (2024).

\bibitem{shakovnotes}
A. Shakov (working notes), \emph{B7-research-integration.md} and related files in the repository \texttt{C:/Users/anton/OneDrive/Desktop/mathAI/notes/}.

\bibitem{mayer}
D. Mayer, \emph{On a $\zeta$-function related to the continued fraction transformation}, Bull. Soc. Math. France \textbf{104} (1976), 195--203; also \emph{The Ruelle-Araki Transfer Operator in Classical Statistical Mechanics}, Lecture Notes in Physics \textbf{123}, Springer (1980).

\bibitem{mayer-roepstorff}
D. Mayer and G. Roepstorff, \emph{On the relaxation time of Gauss's continued-fraction map. I.}, J. Stat. Phys. \textbf{47} (1987), 149--171.

\bibitem{wirsing}
E. Wirsing, \emph{On the theorem of Gauss--Kuzmin--L\'evy and a Frobenius-type theorem for function spaces}, Acta Arith. \textbf{24} (1974), 507--528.

\bibitem{ruelle}
D. Ruelle, \emph{Thermodynamic Formalism}, 2nd ed., Cambridge Mathematical Library (2004).

\bibitem{parry-pollicott}
W. Parry and M. Pollicott, \emph{Zeta Functions and the Periodic Orbit Structure of Hyperbolic Dynamics}, Ast\'erisque \textbf{187--188} (1990).

\bibitem{pollicott-vytnova}
M. Pollicott and P. Vytnova, \emph{Hausdorff dimension estimates applied to Lagrange and Markov spectra, Zaremba theory, and limit sets of Fuchsian groups}, arXiv:1611.09276 (2016--2022).

\bibitem{jenkinson-pollicott}
O. Jenkinson and M. Pollicott, \emph{Computing the dimension of dynamically defined sets: $E_2$ and bounded continued fractions}, Ergodic Theory Dynam. Systems \textbf{21} (2001), 1429--1445.

\bibitem{hensley1992}
D. Hensley, \emph{Continued fraction Cantor sets, Hausdorff dimension, and functional analysis}, J. Number Theory \textbf{40} (1992), 336--358.

\bibitem{hensley-book}
D. Hensley, \emph{Continued Fractions}, World Scientific (2006).

\bibitem{khinchin}
A. Y. Khinchin, \emph{Continued Fractions}, University of Chicago Press (1964).

\bibitem{iwaniec1976}
H. Iwaniec, \emph{The half dimensional sieve}, Acta Arith. \textbf{29} (1976), 69--95.

\bibitem{iwaniec1978}
H. Iwaniec, \emph{Almost-primes represented by quadratic polynomials}, Invent. Math. \textbf{47} (1978), 171--188.

\bibitem{iwaniec-kowalski}
H. Iwaniec and E. Kowalski, \emph{Analytic Number Theory}, AMS Colloquium Publications \textbf{53} (2004).

\bibitem{deshouillers-iwaniec}
J.-M. Deshouillers and H. Iwaniec, \emph{Kloosterman sums and Fourier coefficients of cusp forms}, Invent. Math. \textbf{70} (1982), 219--288.

\bibitem{landau-ramanujan}
E. Landau, \emph{\"Uber die Einteilung der positiven ganzen Zahlen in vier Klassen}, Arch. Math. Phys. (3) \textbf{13} (1908), 305--312.

\bibitem{bourgain-gamburd}
J. Bourgain and A. Gamburd, \emph{Uniform expansion bounds for Cayley graphs of $\SL_2(\F_p)$}, Ann. of Math. (2) \textbf{167} (2008), 625--642.

\bibitem{bourgain-kontorovich}
J. Bourgain and A. Kontorovich, \emph{On Zaremba's conjecture}, Ann. of Math. (2) \textbf{180} (2014), 137--196.

\bibitem{magee-oh-winter}
M. Magee, H. Oh, and D. Winter, \emph{Uniform congruence counting for Schottky semigroups in $\SL_2(\Z)$}, J. Reine Angew. Math. \textbf{753} (2019), 89--135.

\bibitem{furstenberg-kifer}
H. Furstenberg and Y. Kifer, \emph{Random matrix products and measures on projective spaces}, Israel J. Math. \textbf{46} (1983), 12--32.

\bibitem{stern-brocot}
N. Calkin and H. Wilf, \emph{Recounting the rationals}, Amer. Math. Monthly \textbf{107} (2000), 360--363.

\bibitem{selberg-parity}
A. Selberg, \emph{On elementary methods in prime-number theory and their limitations}, Tolfte Skand. Mat. Kongr. (1949), 13--22.

\end{thebibliography}

\end{document}
